\textbf{Selección vigente CEN-41.} Se sustituyen los dos LDO negativos TPS7A3301 por dos positivos adicionales, total cuatro TPS7A4700 en cuatro celdas flotantes. Se seleccionan trim OEM1,72 kohm, cuatro LTC4365 y ocho MOS de corte. Las conexiones y pasivos negativos que siguen pertenecen al antecedente sustituido, sin otra compra. Se conservan cuatro interfaces GCI-38 y cuatro TEN/PTCB; secuencia bipolar, ventana válida, histéresis y respuesta dinámica siguen pendientes. El presupuesto DC global vigente con ENV-40 es356,7552 W; las pérdidas posteriores a TEN no se agregan otra vez.

\section{Posregulación de la medida operacional (CEN/ECI-38)}
\label{sec:sd4-posregulacion-cen38}
Se conservan los cuatro LEM IT 200-S de CEN-35 exclusivamente para diagnóstico. Se seleccionan dos Texas Instruments TPS7A4700RGWR positivos y dos TPS7A3301RGWR negativos: un par por banco, montado en placa propia cerrada. Cada regulador alimenta los dos sensores de su banco; se reserva 0,56 A por polaridad, frente a capacidad de regulación 1 A. Los cuatro TEN 30-2413 se conservan como prerreguladores aislados. Se reemplaza su salida de proyecto 15 V por objetivo ajustado 16,5 V, dentro de ajuste OEM $\pm10\,\%$; nunca se conectan directamente al LEM. Un TEN sin ese ajuste no ofrece margen suficiente para conservar una salida regulada de este diseño \parencites[pp.~1 y 3]{lote35-ten30}[pp.~3--6]{lote38-positive}[pp.~4--6]{lote38-negative}.

Se fija salida positiva nominal 14,9 V mediante ANY-OUT: GND en pines 4/5 (6,4 V cada uno), 10 (0,4 V), 11 (0,2 V) y 12 (0,1 V), añadidos a base 1,4 V; pines 6/8/9 quedan abiertos. SENSE pin 3 toma la salida en la carga; GND pin 7 y pad térmico van a referencia 0, IN pines 15/16 a prerregulación positiva y OUT pines 1/20 a sensores. EN pin 13 sigue IN en esta versión; no se inventa un circuito de secuenciación ya resuelto. Para negativo se fijan R1=118 k$\Omega$ entre OUT/FB y R2=10,1 k$\Omega$ entre FB/0, tolerancia de proyecto 0,1\,\%; $|V_{\rm nom}|=1,175(1+118/10,1)=14,9027$ V. TPS7A3301 comparte números de IN/OUT/GND/EN en encapsulado RGW, pero pin 3 es FB y pin 14 NR/SS. No se copia la programación ANY-OUT al negativo \parencites[pp.~3--4 y 14--16]{lote38-positive}[pp.~4 y 16--17]{lote38-negative}.

Cada LDO lleva entrada 10 $\mu$F, salida con capacidad efectiva mínima 20 $\mu$F y nominal 47 $\mu$F cerámica, y NR/SS de 1 $\mu$F; en negativo se fija CFF 10 nF entre FB y OUT según recomendación OEM. Son valores de placa propios, no un módulo comercial certificado; selección final de condensadores debe conservar capacidad efectiva bajo tensión/temperatura y estabilidad. Se fija trazado Kelvin de 0/SENSE y planos térmicos separados de potencia primaria. La polaridad negativa de los condensadores/cables se identifica expresamente. La referencia 0 de cada banco permanece independiente del otro banco y dentro del límite de modo común a tierra del sensor \parencites[pp.~3--4 y 14--18]{lote38-positive}[pp.~4 y 16--18]{lote38-negative}.

\textbf{Margen calculado.} Ambos reguladores publican exactitud global $\pm2,5\,\%$ dentro de sus condiciones de entrada, carga y temperatura; no se usa 307/325 mV típico como caída máxima. A 1 A los máximos de dropout son 450 mV positivo y 800 mV negativo. Para utilizar exactitud global se conserva además entrada al menos 1 V sobre la salida nominal: se fija aceptación de prerregulación $|V_{\rm pre}|\ge16,0$ V en operación, no sólo una comparación con dropout. El cribado anterior de tolerancia estática 2,2\,\%, aplicado a 16,5 V y con 0,05 V de reserva propia de ajuste, da 16,087--16,913 V; esa aplicación no prueba exactitud/rizado de la variante ajustada. Deben comprobarse en la placa cargada, con entrada 18--36 V del TEN y ambiente previsto. La resistencia exacta de trim y su conexión se definirán desde la nota OEM vigente antes del pedido; no se inventa un valor a partir de una ecuación no verificada \parencites[pp.~3--4]{lote35-ten30}[p.~6]{lote38-positive}[p.~6]{lote38-negative}.

Para positivo la banda global calculada es 14,5275--15,2725 V. Para negativo se añaden extremos del cociente de resistencias 0,1\,\% y $I_{\rm FB}\le100$ nA: $1,175[1+118(1\pm0,001)/(10,1(1\mp0,001))]$, multiplicado por 0,975/1,025, más reserva bilateral de 0,0119 V de corriente FB. Resulta banda de magnitud 14,4915--15,3154 V. Se fija caída máxima de cable 0,10 V y desviación dinámica instalada de magnitud no mayor que 0,10 V, ambas reservas propias que deben demostrarse: banda negativa prevista 14,2915--15,4154 V, dentro de 14,25--15,75 V del LEM. No se suman otra vez regulación de línea/carga incluidas en exactitud global. A 0,56 A y prerregulación 16,913 V, se reserva disipación propia 1,50 W por LDO, 6 W globales; a ambiente 50 $^\circ$C se exige resistencia térmica instalada suficiente para mantener unión por debajo de 100 $^\circ$C. La placa no queda térmicamente verificada por la resistencia térmica típica de una ficha \parencites[p.~5]{lote35-lem200}[pp.~5--6 y 18]{lote38-positive}[pp.~5--6 y 19]{lote38-negative}.

\textbf{Fallo y arranque todavía delimitados.} La posregulación seleccionada establece margen estático y capacidad de corriente; no demuestra que un LDO con corto IN--OUT preserve 15,75 V. Protección térmica/límite de corriente interno tampoco es protección independiente de sobretensión. El TEN mantiene OVP 125\,\% típica, rizado/transitorios típicos y tiempo de arranque típico; no se convierten en máximos del conjunto. Se conserva abierta la protección independiente de cada polaridad y la coordinación ante pérdida/asimetría de referencia, cortos, inversión y arranque. No se admite tensión inferior/superior como estado válido de medida, ni se conecta la placa al primario peligroso hasta acreditar esa protección y los extremos instalados. Las pruebas de alimentación se realizan con circuitos primarios aislados/desenergizados: la ficha LEM prohíbe conectar/desconectar su fuente con primario conectado a partes activas. La primera lectura sólo queda disponible después de ambas alimentaciones válidas, estado LEM ON y tres lecturas frescas coherentes; es criterio de diagnóstico propio, sin tiempo total de seguridad acreditado \parencites[pp.~2--3 y 10--11]{lote35-lem200}[pp.~3--4]{lote35-ten30}.

\section{Interfaz seleccionada de estado LEM a DI24V (GCI-38)}
\label{sec:sd4-interfaz-gci38}
Se seleccionan cuatro interfaces de placa propia, una por sensor, con MOSFET P Nexperia PMV100EPA y zéner BZX84-C12. Se conservan las cuatro DI de WAGO 750-430 reservadas en ENV-34. PMV100EPA admite $|V_{\rm DS}|=60$ V, $|V_{\rm GS}|=20$ V y publica $R_{\rm DS,on}\le0,276~\Omega$ a unión 175 $^\circ$C y excitación $-10$ V; pines 1=G, 2=S y 3=D. El zéner publica 11,4--12,7 V a 5 mA. No se atribuye nivel SIL, aislamiento primario ni ensayo de módulo industrial a estas placas \parencites[pp.~1--3 y 6]{lote38-pmos}[pp.~2 y 6]{lote38-zener}.

En cada canal S toma +24 V local; G tiene resistencia 10 k$\Omega$/5\,\% a S y zéner con cátodo a S/ánodo a G. G conecta a colector LEM pin 8 mediante 1,2 k$\Omega$/5\,\%, mínimo 0,5 W; emisor pin 3 toma 0 V de campo local. D entrega DI a través de Vishay AC03 430 $\Omega$/5\,\%, referencia AC03000004300JAC00. Una resistencia de 10 k$\Omega$/5\,\% entre DI y 0 se coloca al extremo de entrada WAGO, después del cable de señal. El opto ON tira de G, enciende el P-MOS y DI recibe alto; opto OFF deja G en S y el pull-down fuerza bajo. No se unen pines 1/4 de medida LEM a emisor de estado por conveniencia: pantalla, retorno analógico y retorno de campo siguen sus circuitos definidos \parencites[pp.~8 y 10]{lote35-lem200}[p.~2]{lote38-pmos}[pp.~1--2]{lote38-ac}.

Se fija ventana de campo 21,6--28,8 V durante operación y extremo de cálculo de falla 30 V. Con $V_Z=11,4$--12,7 V, $V_{\rm CE}\le0,2$ V y extremos de resistencia 1,2 k$\Omega$, corriente de opto queda entre 6,90 y 16,32 mA, dentro de 2--30 mA LEM. El consumo de la resistencia G--S se descuenta al verificar corriente de zéner: el mínimo calculado es 5,57 mA. Se comprueba temperatura/dinámica del zéner en placa; valores a 5 mA no se presentan como clamp idéntico a cualquier corriente. La resistencia serie disipa hasta 0,31 W; se fija 0,5 W mínimo. La AC03 publica 3 W a 40 $^\circ$C y 2,5 W a 70 $^\circ$C: a salida en corto y 30 V, $I\le30/(430\cdot0,95)=73,44$ mA y $P_R\le2,203$ W. Debe conservar separación térmica de LEM/cables. Ese cálculo protege el circuito seleccionado bajo sus condiciones; no detecta por sí solo un corto que produce DI alto \parencites[p.~8]{lote35-lem200}[p.~1]{lote38-ac}[pp.~3 y 6]{lote38-pmos}.

WAGO publica bajo $-3$--$+5$ V, alto 15--30 V y corriente de entrada 2,8 mA típica. Se fija límite propio de aceptación de corriente DI de 6 mA; no se lo presenta como máximo OEM. Con extremo bajo 21,6 V, AC03 máxima 451,5 $\Omega$, pull-down mínimo 9,5 k$\Omega$ y resistencia MOS 0,276 $\Omega$, resulta $V_{\rm DI}\ge[21,6-0,006(451,5+0,276)]/[1+(451,5+0,276)/9500]=18,0318$ V, mayor que 15 V. Si corriente instalada supera 6 mA o cambia la ventana de campo, se recalcula y no se libera el canal. En estado OFF se exige tensión medida $\le5$ V y comprobación de fugas a 50 $^\circ$C; no se extrapola la fuga MOS publicada a 25 $^\circ$C como máximo completo de todo el circuito. El filtro WAGO de 3 ms no se transforma en latencia máxima conjunta \parencite[p.~14]{lote38-wagodi}.

Un cable cortado en colector/emisor abre el mando y un corte de señal hacia DI deja su pull-down local: lectura baja bajo alimentación/retorno íntegros. Pérdida de retorno común se trata como conjunto inválido, sin afirmar que todos los canales forzosamente caerán a cero. Corto colector--emisor, G forzado a bajo, MOS S--D en corto o cable a +24 V pueden simular estado sano; se registran como fallos no cubiertos por esta interfaz simple. La aceptación por canal inyecta ON/OFF, corta cada conductor, prueba corto de señal y comprueba estados con primaria aislada; también verifica señal congelada/edad y retiro de validez en aplicación. No se conecta este estado a permiso XPS ni se gobierna recarga. La disponibilidad de diagnóstico no habilita potencia y la alarma de 17 A de CEN-35 conserva su carácter operacional.

\textbf{Conciliación.} Se agregan cuatro LDO, cuatro MOSFET, cuatro zéner, cuatro AC03 y pasivos de cuatro interfaces; son dos placas de banco, sin nuevas DI ni nuevos TEN/PTCB/QUINT. LDO y pasivos de posregulación permanecen aguas abajo de los cuatro TEN ya reservados: no se agregan 6 W de pérdida como segunda carga DC/AC. Se reserva adicionalmente 1 W de campo por interfaz, 4 W globales/2 W por QUINT, excluyendo la DI ya reservada en ENV-34. Sobre base integrada 348,7552 W, el subtotal reservado pasa a 352,7552 W DC, 176,3776 W por QUINT antes de receptores pendientes. La reserva propia de 192 W para cuatro TEN se conserva mientras se verifica corriente/arranque reales a ajuste16,5 V; no se deriva una eficiencia mínima del 91\,\% típico a salida nominal15 V. Falla/coord, capacidad UPS y grupo, autonomía, recarga OEM20 A y FP/PUE anual siguen abiertos. RT-06.07/.08/.11 no se completan con esta selección parcial.
